\documentclass[12pt,a4paper]{article}

\usepackage[utf8]{inputenc}
\usepackage[T2A]{fontenc}
\usepackage[russian]{babel}
\usepackage[margin=2cm]{geometry}
\usepackage{amsmath}
\usepackage{listings}
\usepackage{xcolor}
\usepackage{hyperref}
\usepackage{booktabs}
\usepackage{indentfirst}

\setlength{\parindent}{1.25cm}

\definecolor{codebg}{rgb}{0.95,0.95,0.95}
\definecolor{codegreen}{rgb}{0.0,0.5,0.0}
\definecolor{codegray}{rgb}{0.5,0.5,0.5}
\definecolor{codepurple}{rgb}{0.58,0,0.82}

\lstdefinestyle{javastyle}{
    backgroundcolor=\color{codebg},
    commentstyle=\color{codegreen},
    keywordstyle=\color{blue},
    stringstyle=\color{codepurple},
    numberstyle=\tiny\color{codegray},
    basicstyle=\ttfamily\footnotesize,
    breakatwhitespace=false,
    breaklines=true,
    captionpos=b,
    keepspaces=true,
    numbers=left,
    numbersep=5pt,
    showspaces=false,
    showstringspaces=false,
    showtabs=false,
    tabsize=2,
    language=Java,
    frame=single
}
\lstset{style=javastyle}

\begin{document}

\begin{titlepage}
    \setcounter{page}{1}
    \thispagestyle{empty}

    \begin{center}
        { \large МИНИСТЕРСТВО НАУКИ И ВЫСШЕГО ОБРАЗОВАНИЯ \\
            РОССИЙСКОЙ ФЕДЕРАЦИИ }

        { «Санкт-Петербургский политехнический университет Петра Великого» }

        { Институт компьютерных наук и кибербезопасности \\
            Высшая школа технологий искусственного интеллекта \\
            Направление: 02.03.01. «Математика и компьютерные науки» \\ }

        \vspace{3cm}
        { \large Отчет о выполнении лабораторной работы \\
            По дисциплине: «Архитектура суперкомпьютеров» \\
            На тему: «Реализация алгоритма перевода чисел формата BF16 из двоичной системы счисления в десятичную и обратно» }

    \end{center}

    \vspace{3cm}

    \begin{center}
        \noindent
        \begin{minipage}[t]{\textwidth}
            \raggedright Студент  \\ группы 5130201/40101 \hfill
            \underline{\hspace{3cm}} Храмов В.Д.
        \end{minipage}

        \vspace{1cm}

        \noindent
        \begin{minipage}[t]{\textwidth}
            \raggedright Преподаватель \hfill
            \underline{\hspace{3cm}}  Чуватов М.В.
        \end{minipage}
    \end{center}

    \vspace{1cm}

    \begin{flushright}
        \large
        «\underline{\hspace{1cm}}» \underline{\hspace{3cm}} 2026 г.
    \end{flushright}

    \vspace{0.5cm}

    \vfill

    \begin{center}
        Санкт-Петербург \\
        2026
    \end{center}

\end{titlepage}

\section{Цель работы}

Изучить формат представления чисел с плавающей точкой \texttt{bfloat16} (BF16),
разработать и программно реализовать алгоритмы:
\begin{itemize}
    \item перевода десятичного числа в его 16-битное представление BF16;
    \item обратного перевода 16-битного представления BF16 в десятичное число.
\end{itemize}

\section{Теоретическая часть}

\subsection{Формат BF16}

\texttt{bfloat16} — это 16-битный формат числа с плавающей точкой,
использующий ту же ширину поля экспоненты, что и стандартный 32-битный
формат \texttt{float} (IEEE~754 single precision), но с урезанной мантиссой.
Формат состоит из трёх полей:

\begin{center}
\begin{tabular}{|c|c|c|}
\hline
Знак (S) & Экспонента (E) & Мантисса (M) \\
\hline
1 бит & 8 бит & 7 бит \\
\hline
\end{tabular}
\end{center}

Смещение (bias) экспоненты равно $127$. Значение нормализованного числа
вычисляется по формуле:

\begin{equation}
value = (-1)^S \times 1.M \times 2^{(E - 127)}
\end{equation}

где $1.M$ — мантисса с неявной единицей впереди.

\subsection{Специальные случаи}

\begin{itemize}
    \item $E = 0,\ M = 0$ — число равно $\pm 0$;
    \item $E = 0,\ M \neq 0$ — денормализованное число: $value = (-1)^S \times 0.M \times 2^{-126}$;
    \item $E = 255,\ M = 0$ — $\pm\infty$;
    \item $E = 255,\ M \neq 0$ — NaN.
\end{itemize}

\subsection{Диапазон и точность}

Поскольку под экспоненту выделено столько же бит (8), сколько и в
обычном \texttt{float32}, диапазон представимых значений BF16 совпадает
с диапазоном \texttt{float32} и составляет примерно от $1.18 \times 10^{-38}$
до $3.39 \times 10^{38}$. Плата за компактность формата — точность:
из-за короткой (7-битной) мантиссы BF16 сохраняет лишь около 2--3
десятичных знаков точности, в то время как \texttt{float32} сохраняет
около 7.

\section{Алгоритм перевода}

\subsection{Из десятичной системы в BF16}

\begin{enumerate}
    \item Определить знак числа $S$ и перейти к работе с его модулем.
    \item Нормализовать число к виду $m \times 2^{e}$, где $1.0 \le m < 2.0$,
    последовательным делением или умножением на 2.
    \item Дробную часть нормализованной мантиссы $m - 1$ перевести
    в двоичный вид методом умножения на 2 с фиксированной точностью
    (\textit{с запасом}, чтобы было по чему округлять).
    \item Округлить полученную мантиссу до 7 бит по правилу
    \textit{round-to-nearest-even} (округление к чётному в случае
    "ровно посередине"), учитывая перенос в экспоненту при переполнении
    мантиссы.
    \item Вычислить смещённую экспоненту $E = e + 127$ и проверить границы
    диапазона (переполнение $\to \infty$, антипереполнение $\to 0$).
    \item Собрать итоговые 16 бит: $S \,|\, E \,|\, M$.
\end{enumerate}

\subsection{Из BF16 в десятичную систему}

\begin{enumerate}
    \item Разбить 16 бит на поля $S$ (1 бит), $E$ (8 бит), $M$ (7 бит).
    \item Обработать специальные случаи ($\pm 0$, $\pm\infty$, NaN,
    денормализованные числа).
    \item Перевести мантиссу $M$ обратно в дробное десятичное число:
    $\sum_{i=1}^{7} bit_i \times 2^{-i}$.
    \item Вычислить итоговое значение по формуле
    $(-1)^S \times (1 + M) \times 2^{(E - 127)}$.
\end{enumerate}

\section{Реализация}

\subsection{Перевод целой и дробной части в двоичный вид}

Перевод целой части выполняется классическим методом последовательного
деления на 2 с накоплением остатков; перевод дробной части — методом
последовательного умножения на 2:

\begin{lstlisting}[caption={Перевод целой части в двоичный вид}]
private static String decimalToBinaryInteger(int num) {
    StringBuilder binary = new StringBuilder();

    if (num == 0) {
        return "0";
    }

    while (num > 0) {
        int remainder = num % 2;
        num /= 2;
        binary.append(remainder);
    }

    return binary.reverse().toString();
}
\end{lstlisting}

\begin{lstlisting}[caption={Перевод дробной части в двоичный вид}]
private static String decimalToBinaryFraction(double frac, int precision) {
    StringBuilder binary = new StringBuilder();
    int countPrecision = 0;

    while (frac != 0 && countPrecision != precision) {
        frac *= 2;
        if (frac >= 1) {
            binary.append(1);
            frac -= 1;
        } else {
            binary.append(0);
        }
        countPrecision++;
    }

    while (countPrecision != precision) {
        binary.append(0);
        countPrecision++;
    }
    return binary.toString();
}
\end{lstlisting}

\subsection{Нормализация}

Нормализация объединяет целую и дробную части двоичного
числа в одну строку \texttt{combined}. Позиция двоичной точки определяется
по длине целой части и сохраняется в переменной \texttt{pointPosition}.

С помощью метода \texttt{indexOf('1')} находится позиция первой единицы в
объединённой строке. Если единица отсутствует, число считается равным нулю,
и выбрасывается исключение \texttt{IllegalArgumentException}.

Показатель степени вычисляется как разница между позицией двоичной точки и
позицией первой единицы:
\[
exponent = pointPosition - firstOne - 1.
\]

Мантисса формируется из последовательности битов, расположенных после первой
единицы, с помощью метода \texttt{substring(firstOne + 1)}. Таким образом,
метод возвращает показатель степени и мантиссу нормализованного двоичного
числа.

\begin{lstlisting}[caption={Нормализация через деление/умножение на 2}]
private static Normalized normalize(String binary, String binaryFraction) {
    String combined = binary + binaryFraction;
    int pointPosition = binary.length();

    int firstOne = combined.indexOf('1');
    if (firstOne == -1) {
        throw new IllegalArgumentException("");
    }

    int exponent = pointPosition - firstOne - 1;
    String mantissa = combined.substring(firstOne + 1);

    return new Normalized(exponent, mantissa);
}
\end{lstlisting}

\subsection{Округление мантиссы (round-to-nearest-even)}

Так как исходная точность мантиссы (32 бита) значительно превышает
целевые 7 бит BF16, отбрасываемые биты необходимо корректно округлять.
Используется классическое правило: если бит-\textit{guard} (8-й бит) равен
0 — округление не требуется; если он равен 1 и после него есть хотя бы
одна единица — округление вверх; если это граничный случай "ровно
посередине" — округление производится к чётному значению последнего
сохраняемого бита. При переполнении мантиссы (например, \texttt{1111111}
округляется до \texttt{0000000}) единица переноса добавляется к
экспоненте.

\begin{lstlisting}[caption={Округление мантиссы до 7 бит}]
private static Rounded roundMantissa(String mantissa, int exponent, int targetBits) {
    StringBuilder padded = new StringBuilder(mantissa);
    while (padded.length() < targetBits + 1) {
        padded.append('0');
    }

    String kept = padded.substring(0, targetBits);
    char guardBit = padded.charAt(targetBits);
    String remainder = padded.length() > targetBits + 1 ? padded.substring(targetBits + 1) : "";

    boolean roundUp;
    if (guardBit == '0') {
        roundUp = false;
    } else if (remainder.contains("1")) {
        roundUp = true;
    } else {
        char lastKeptBit = kept.charAt(kept.length() - 1);
        roundUp = (lastKeptBit == '1');
    }

    int keptValue = kept.isEmpty() ? 0 : Integer.parseInt(kept, 2);
    if (roundUp) {
        keptValue += 1;
    }

    int maxValue = (1 << targetBits) - 1;
    if (keptValue > maxValue) {
        keptValue = 0;
        exponent += 1;
    }

    String mantissaBits = toFixedLengthBinary(keptValue, targetBits);
    return new Rounded(exponent, mantissaBits);
}
\end{lstlisting}

\subsection{Кодирование десятичного числа в BF16}

Кодирование выполняется в несколько этапов. Сначала определяется знак
числа и выделяется его абсолютное значение. Отдельно обрабатываются
специальные значения: \texttt{NaN}, положительная и отрицательная
бесконечность, а также ноль.

Для ненулевого конечного числа целая и дробная части преобразуются в
двоичное представление. Для дробной части используется фиксированная
точность в 32 двоичных разряда. Полученные части передаются в метод
\texttt{normalize}, который определяет показатель степени и мантиссу
нормализованного числа.

После нормализации мантисса округляется до 7 бит с помощью метода
\texttt{roundMantissa}. К полученному показателю степени прибавляется
смещение 127, используемое в формате BF16. Если после применения
смещения показатель выходит за допустимый диапазон, значение
обрабатывается как переполнение или антипереполнение.

Затем знак, показатель степени и мантисса объединяются в одно 16-битное
значение с помощью побитовых операций. Знак занимает старший бит,
показатель степени --- следующие 8 бит, а мантисса --- младшие 7 бит.

\begin{lstlisting}[caption={Кодирование double в BF16}]
public static int encode(double value) {
    int sign = (value < 0) ? 1 : 0;
    double absValue = Math.abs(value);

    if (Double.isNaN(value)) {
        return (sign << 15) | (0xFF << 7) | 0x40;
    }
    if (Double.isInfinite(value)) {
        return (sign << 15) | (0xFF << 7);
    }
    if (absValue == 0.0) {
        return sign << 15;
    }

    int intPart = (int) absValue;
    double fracPart = absValue - intPart;

    String binaryInteger = decimalToBinaryInteger(intPart);
    String binaryFraction = decimalToBinaryFraction(fracPart, 32);

    Normalized norm = normalize(binaryInteger, binaryFraction);
    Rounded rounded =
        roundMantissa(norm.mantissa(), norm.exponent(), 7);

    int biasedExponent = rounded.exponent() + 127;

    if (biasedExponent >= 255) {
        return (sign << 15) | (0xFF << 7);
    }
    if (biasedExponent <= 0) {
        return sign << 15;
    }

    int mantissaBits = Integer.parseInt(rounded.mantissa(), 2);

    return (sign << 15) | (biasedExponent << 7) | mantissaBits;
}
\end{lstlisting}

\subsection{Декодирование BF16 в десятичное число}

Обратный перевод восстанавливает значение непосредственно по формуле
стандарта, без использования встроенных средств Java для работы с битами
\texttt{float}:

\begin{lstlisting}[caption={Декодирование BF16 в double}]
public static double decode(int bits) {
    bits = bits & 0xFFFF;

    int sign = (bits >> 15) & 0x1;
    int exponent = (bits >> 7) & 0xFF;
    int mantissa = bits & 0x7F;

    double signMultiplier = (sign == 0) ? 1.0 : -1.0;

    if (exponent == 0 && mantissa == 0) {
        return signMultiplier * 0.0;
    }
    if (exponent == 255) {
        return (mantissa == 0)
                ? signMultiplier * Double.POSITIVE_INFINITY
                : Double.NaN;
    }
    if (exponent == 0) {
        double mantissaValue = mantissaToDecimal(mantissa, 7);
        return signMultiplier * mantissaValue * Math.pow(2, -126);
    }

    double mantissaValue = mantissaToDecimal(mantissa, 7);
    int realExponent = exponent - 127;

    return signMultiplier * (1 + mantissaValue) * Math.pow(2, realExponent);
}
\end{lstlisting}

\section{Тестирование}

Корректность реализации проверена набором модульных тестов,
охватывающих: точно представимые значения, ноль, отрицательные числа,
потерю точности, округление к чётному в граничных случаях, специальные
значения (NaN, $\pm\infty$) и выход за границы диапазона (переполнение и
антипереполнение). Результаты работы программы на контрольных значениях
приведены в таблице~\ref{tab:results} (формат поля битов: S EEEEEEEE MMMMMMM).

\begin{table}[h]
\centering
\caption{Результаты кодирования и декодирования чисел в формате BF16}
\label{tab:results}
\small
\begin{tabular}{@{}c c c c c p{3.0cm}@{}}
\toprule
Исходное значение &
Знак &
Экспонента &
Мантисса &
Биты BF16 &
Декодированное значение \\
\midrule

$3.0$
& $0$
& $10000000$
& $1000000$
& $0\ 10000000\ 1000000$
& $3.0$ \\

$5.25$
& $0$
& $10000001$
& $0101000$
& $0\ 10000001\ 0101000$
& $5.25$ \\

$0.1$
& $0$
& $01111011$
& $1001101$
& $0\ 01111011\ 1001101$
& $0.10009765625$ \\

$-2.5$
& $1$
& $10000000$
& $0100000$
& $1\ 10000000\ 0100000$
& $-2.5$ \\

$100.0$
& $0$
& $10000101$
& $1001000$
& $0\ 10000101\ 1001000$
& $100.0$ \\

$0.0$
& $0$
& $00000000$
& $0000000$
& $0\ 00000000\ 0000000$
& $0.0$ \\

$65536.0$
& $0$
& $10001111$
& $0000000$
& $0\ 10001111\ 0000000$
& $65536.0$ \\

$10^{40}$
& $0$
& $11111111$
& $0000000$
& $0\ 11111111\ 0000000$
& $+\infty$ \\

$10^{-40}$
& $0$
& $00000000$
& $0000000$
& $0\ 00000000\ 0000000$
& $0.0$ \\

\bottomrule
\end{tabular}
\end{table}

Значение $0.1$ демонстрирует ожидаемую потерю точности: из-за короткой
7-битной мантиссы BF16 не может точно представить это число, и после
декодирования получается близкое, но не идентичное значение
($0.10009765625$). Значения $3.0$, $5.25$, $100.0$ и $65536.0$, напротив,
представимы в BF16 точно, что подтверждает корректность работы алгоритма
нормализации и округления.

\section{Выводы}

В ходе выполнения лабораторной работы был изучен формат чисел с плавающей
точкой BF16 и самостоятельно реализован полный цикл перевода чисел между
десятичной и двоичной системами счисления по этому формату — без
использования готовых средств языка Java для работы с битовым
представлением \texttt{float}.

\end{document}